\documentclass[11pt]{article}
\usepackage[margin=1in]{geometry}
\usepackage{amsmath,amssymb}
\usepackage{graphicx}
\usepackage{booktabs}
\usepackage{xcolor}
\usepackage{hyperref}

\newcommand{\E}{\mathbb{E}}
\newcommand{\Norm}{\mathcal{N}}
\newcommand{\code}[1]{\texttt{#1}}

\title{Analysis plots for the joint disturbance / process-error filter\\
       (\code{ORJointUpdate} in \code{OR\_joint\_update2.Rmd})}
\author{Claude}
\date{October 7, 2026}

\emergencystretch=3em
\begin{document}
\maketitle

This note explains the diagnostic plots that \code{OR\_joint\_update\_analysis.Rmd}
makes from runs of \code{ORJointUpdate}, and gives the formula for every quantity they
show:
\begin{enumerate}
  \item filter divergence at the final observation of each site, with an alternative
        statistic (the predictive PIT) aimed at the same question;
  \item the probability of detecting a disturbance in 2004 against the observed biomass
        loss from 2003 to 2004.
\end{enumerate}
Section~\ref{sec:model} defines the filter's notation and Section~\ref{sec:pred} the
one-step-ahead predictive. Sections~\ref{sec:div} and~\ref{sec:det} define the plotted
quantities. Section~\ref{sec:runs} describes the experiment grid and
the 5-site test run in \code{experiments/testRun}, and shows its plots.

%================================================================================
\section{The filter}
\label{sec:model}

\subsection{Data, time axis and state}

\begin{itemize}
  \item \textbf{Sites and years.} $s$ indexes sites (pixels of tile h03v04). The filter runs
  on annual timesteps $t = 1,\dots,T$, $T = 28$, where timestep $t$ is calendar year
  $1989 + t$ (1990--2017; \code{atime}). The biomass product covers 1990--2017 and the
  disturbance product covers 1985--2012. Both are looked up by calendar year, so the
  product's flag for year $Y$ is applied at the timestep for year $Y$.
  \item \textbf{State.} $x_t = (x_{t,1}, x_{t,2}, x_{t,3})^\top$ is the vector of three carbon
  pools (kg\,C\,m$^{-2}$) of VSEM, the Very Simple Ecosystem Model in the R package
  \code{BayesianTools}. The code treats them as leaf, soil and wood. The wood pool
  $x_{t,3}$ is the one compared with above-ground biomass. One VSEM year of the state
  $x$, driven by that site's daily photosynthetically active radiation (PAR, from
  Daymet), is written $f_t(x)$.
  \item \textbf{Observation.} $y_t$ is the Landsat-based above-ground biomass (AGB) of the
  site in year $1989+t$. It is converted from Mg\,ha$^{-1}$ to kg\,C\,m$^{-2}$ by
  $\times 0.5$ (biomass to carbon) $\times 0.1$. The observation model is
  \[
    y_t \mid x_t \sim \Norm(Hx_t,\ R_t), \qquad H = (0,0,1), \qquad R_t = m_R\,y_t + b_R ,
  \]
  where $R_t$ is the observation-error variance and $(m_R, b_R)$ is \code{Rparams}.
  ``Heteroskedastic'' (\code{het}) runs have $m_R > 0$, so $R_t$ grows with biomass.
  ``Constant'' (\code{con}) runs have $m_R = 0$.
\end{itemize}

\subsection{Disturbance labels and their prior}

Each year the site carries one label $k \in \{0, 1, \dots, D\}$:
\begin{itemize}
  \item $k = 0$ means undisturbed.
  \item $k \ge 1$ is a disturbance sub-class. Each disturbance type of the product
  (other, cut, fire, biotic) is split into the clusters of a Gaussian mixture fitted to
  that type's fraction of wood retained after disturbance. There are $4 + 6 + 3 + 3$
  clusters, so $D = 16$.
  \item Under label $k \ge 1$, the disturbance step $d_k$
  (\code{disturbance.p}) multiplies leaf and wood by random retention fractions. These are
  drawn from $\Norm\big((0.2, r_k),\ \mathrm{diag}(0.04, v_k)\big)$, where $r_k$ and $v_k$
  are the cluster's mean and variance, truncated to $[0, 1]$. Of the carbon removed, 5\%
  goes to the soil pool.
\end{itemize}

The prior probability of label $k$ in year $t$ is $\beta_{k,t}$ (\code{dist\$p[k+1, t]}).
It is built as follows:
\begin{itemize}
  \item \textbf{Background.} Each type gets its background rate (\code{d.prior}). That rate
  is split across the type's clusters in proportion to the cluster weights.
  \item \textbf{With the disturbance product.} In a year the product flags with a type, that
  type's probability is replaced by the product's \emph{user's accuracy} for the type
  (Kennedy et al.\ 2015): the fraction of the product's flags of that type that are
  correct (0.77 other, 0.92 cut, 0.78 fire, 0.78 biotic). It is then split the same way.
  \item \textbf{Undisturbed.} $\beta_{0,t} = 1 - \sum_{k\ge1}\beta_{k,t}$.
\end{itemize}

\subsection{Process error and the precision grid}

Process error is added in the undisturbed label only. Its covariance is
\[
  Q(\tau) = \mathrm{diag}\big(q_1^2,\ q_2^2,\ \tau^{-1}\big),
\]
where $q_1, q_2$ are fixed leaf and soil standard deviations from the calibration
(\code{fit.RData}). $\tau$ is the unknown \emph{precision} (inverse variance) of the wood
process error.

The filter places $\tau$ on a grid $\tau_1 < \dots < \tau_G$ of $G = 200$ points,
spaced evenly in $\log\tau$ from $0.1$ to $10^4$. The prior is Gamma(shape $0.1$,
rate $2\times10^{-4}$). $\pi_0(g)$ is the prior mass of the grid cell around $\tau_g$,
with cell edges at the midpoints between neighbouring grid points, renormalised over the
grid. $\pi_t(g)$ denotes the posterior probability of $\tau_g$ given $y_{1:t}$, the
observations up to year $t$.

\subsection{Posterior representation}

Conditional on $\tau_g$, the posterior of $x_t$ is a mixture of two Gaussians,
$j \in \{1, 2\}$, with:
\begin{itemize}
  \item weights $w_t(g,j)$, where $w_t(g,1) + w_t(g,2) = 1$;
  \item means $m_t(g,j)$ and covariances $P_t(g,j)$.
\end{itemize}
Component $j = 1$ collects the undisturbed label and $j = 2$ the disturbance labels (see
the projection below). At $t = 0$ every component is $\Norm(\bar x_0, P_0)$ with weight
$\tfrac12$, where $\bar x_0, P_0$ are the mean and covariance of an initial-condition
ensemble. That ensemble pairs every site's 1990 AGB with leaf and soil pools from a VSEM
regrowth run of matching wood biomass. It is therefore a cross-site distribution, not
specific to the site being filtered.

\subsection{Forecast step}

\begin{enumerate}
  \item \textbf{Proposal.} Draw $N = 100$ samples $\xi^{(1)},\dots,\xi^{(N)}$ from the
  equal-weight mixture of all $2G$ components, rejecting samples with a non-positive pool.
  The proposal density is
  $q(\xi) = \frac{1}{2G}\sum_{g,j}\Norm\big(\xi;\ m_{t-1}(g,j), P_{t-1}(g,j)\big)$.
  \item \textbf{Propagation.} Push each sample through one VSEM year: $f_t(\xi^{(n)})$. For
  label $k \ge 1$, also apply the disturbance step, giving $d_k(f_t(\xi^{(n)}))$. For
  $k = 0$ the result is left unchanged. Write $z^{(n)}_k$ for the result.
  \item \textbf{Importance weights.} Component $(g,j)$ reweights the shared samples by
  \[
    W_{gj}(n) \propto \frac{\Norm\big(\xi^{(n)};\ m_{t-1}(g,j), P_{t-1}(g,j)\big)}{q(\xi^{(n)})},
    \qquad \sum_n W_{gj}(n) = 1 .
  \]
  \item \textbf{Forecast moments.} The forecast moments of component $(g,j)$ under label
  $k$ are
  \[
    \mu^F_{gjk} = \sum_n W_{gj}(n)\, z^{(n)}_k, \qquad
    \Sigma^F_{gjk} = \frac{\sum_n W_{gj}(n)\,(z^{(n)}_k - \mu^F_{gjk})(z^{(n)}_k - \mu^F_{gjk})^\top}
                         {1 - \sum_n W_{gj}(n)^2}
                   + \mathbf{1}\{k = 0\}\, Q(\tau_g).
  \]
  The denominator is the usual bias correction for a weighted sample covariance. We write
  \[
    C_{gjk} = \big[\Sigma^F_{gjk} - \mathbf{1}\{k=0\}Q(\tau_g)\big]_{33}
  \]
  for the ensemble (sample) variance of the wood pool. The wood forecast variance is then
  $\Sigma^F_{gjk,33} = C_{gjk} + \mathbf{1}\{k = 0\}\,\tau_g^{-1}$.
  \item \textbf{Degenerate components.} If $1 - \sum_n W_{gj}(n)^2 < 10^{-8}$, component
  $(g,j)$'s weights sit on essentially one sample, and its covariance cannot be estimated.
  This happens to components that are far from every proposal sample. Such a component
  is given zero likelihood in the analysis, so zero weight. Its moments are set to the
  unweighted sample moments as a finite placeholder. The forecast weight removed is
  stored per step as \code{hist\$wDrop}. In the test run it occurred once (site 212,
  2005, with no product and GEDI heteroskedastic $R$). The weight removed there was
  $10^{-119}$.
\end{enumerate}

\subsection{Analysis step}

For every $(g,j,k)$, a Kalman update with observation $y_t$ gives the analysis moments
and the component likelihood
\[
  \ell_{gjk} = \Norm\big(y_t;\ \mu^F_{gjk,3},\ \Sigma^F_{gjk,33} + R_t\big).
\]
The marginal likelihood of $y_t$ given $\tau_g$, the $\tau$ posterior, and the analysis
weight of each $(j,k)$ given $\tau_g$ are
\[
  Z_t(g) = \sum_{j,k} w_{t-1}(g,j)\,\beta_{k,t}\,\ell_{gjk}, \qquad
  \pi_t(g) \propto \pi_{t-1}(g)\,Z_t(g), \qquad
  w^A_t(g,j,k) = \frac{w_{t-1}(g,j)\,\beta_{k,t}\,\ell_{gjk}}{Z_t(g)} .
\]
\textbf{Projection.} For each $g$, the $2(D+1)$ analysis components are moment-matched
down to two Gaussians:
\begin{itemize}
  \item $j = 1$ from all components with $k = 0$, with weight
        $w_t(g,1) = \sum_j w^A_t(g,j,0)$;
  \item $j = 2$ from all components with $k \ge 1$, with weight $w_t(g,2) = 1 - w_t(g,1)$.
\end{itemize}
In a year without an observation, the analysis equals the forecast and $\ell_{gjk} = 1$.

\paragraph{Analysis summaries.} The filter stores, for each year, summaries of the
marginal posterior of the wood pool given $y_{1:t}$. That posterior is the Gaussian
mixture
\begin{equation}
  p(x_{t,3} \mid y_{1:t}) = \sum_{g,j,k} \pi_t(g)\, w^A_t(g,j,k)\;
     \Norm\big(x_{t,3};\ \mu^A_{gjk,3},\ \Sigma^A_{gjk,33}\big),
  \label{eq:post}
\end{equation}
where $\mu^A_{gjk}$ and $\Sigma^A_{gjk}$ are the Kalman analysis mean and covariance of
component $(g,j,k)$. The stored summaries are:
\begin{itemize}
  \item the analysis mean $\hat x_t = \E[x_{t,3} \mid y_{1:t}]$ (\code{hist\$mean});
  \item the analysis variance (\code{hist\$var});
  \item the 95\% analysis interval $[L_t, U_t]$ (\code{hist\$ci95}). $L_t$ and $U_t$ are
  the 2.5\% and 97.5\% quantiles of the mixture~\eqref{eq:post}, found by root-finding
  on its cumulative distribution function. Because the mixture can be skewed (for example
  in years with a high disturbance prior), the interval need not be symmetric about
  $\hat x_t$.
\end{itemize}

%================================================================================
\section{The one-step-ahead predictive}
\label{sec:pred}

The second statistic in Plot~1 uses the distribution of $y_t$ given $y_{1:t-1}$, i.e.\
\emph{before} $y_t$ is assimilated. Under the filter above, this predictive is the mixture
\begin{equation}
  p(y_t \mid y_{1:t-1}) = \sum_{g=1}^{G}\sum_{j=1}^{2}\sum_{k=0}^{D}
    \omega_{gjk}\;\Norm\!\big(y_t;\ \mu^F_{gjk,3},\ \Sigma^F_{gjk,33} + R_t\big),
  \qquad
  \omega_{gjk} = \pi_{t-1}(g)\, w_{t-1}(g,j)\, \beta_{k,t}.
  \label{eq:pred}
\end{equation}
\begin{itemize}
  \item \textbf{Integration over process error.} The sum over $g$ uses $\pi_{t-1}$, the
  posterior of the process precision from the previous step. So the predictive averages
  over the process-error variance $1/\tau$ instead of fixing it.
  \item \textbf{Disturbance labels.} The sum over $k$ uses the prior $\beta_{k,t}$, because
  the year-$t$ label is not known before $y_t$ is seen.
  \item \textbf{Weights.} The weights $\omega_{gjk}$ sum to one. Degenerate components get
  $\omega = 0$, and the rest are renormalised.
\end{itemize}

%================================================================================
\section{Plot 1: filter divergence and predictive PIT}
\label{sec:div}

Both statistics ask whether the filter still tracks the data at the end of the series,
and both scale the error by the filter's own uncertainty. Let $T_s$ be the last year with
an observation at site $s$ (2017 for every site here).

\subsection{Filter divergence}

\paragraph{Formula.} The final observation minus the final analysis mean, divided by half
the width of the 95\% analysis interval:
\[
  \delta_s = \frac{y_{T_s} - \hat x_{T_s}}{(U_{T_s} - L_{T_s})/2}.
\]
If the interval is symmetric about $\hat x_{T_s}$, $\delta_s = \pm1$ exactly at its
edges, and $|\delta_s| \le 1$ means the final observation falls inside it.

\paragraph{How to read it.} A calibrated filter should give $|\delta_s| \le 1$ at most
sites. A filter that has \emph{diverged} has a confident analysis that sits away from
the data, giving $|\delta_s| > 1$. The sign shows the direction: $\delta_s > 0$ means the
analysis is below the observation. Two caveats:
\begin{itemize}
  \item $\hat x_{T_s}$ has already assimilated $y_{T_s}$, so it is pulled toward the
  observation. That pull shrinks $\delta_s$, so the statistic flags only strong
  divergence.
  \item $[L_t, U_t]$ is an interval for the latent wood pool. It does not include the
  observation error $R_t$, which acts in the other direction.
\end{itemize}
So $[-1, 1]$ is a rule of thumb, not an exact coverage level.

\paragraph{Panels.} Page~1 of \code{exps\_summary.pdf} plots $\delta_s$ against
$y_{T_s}$, one point per site and one panel per experiment. The band $[-1, 1]$ is shaded,
and the panel title gives the share of sites inside it.

\subsection{Predictive PIT (alternative statistic)}

\paragraph{Formula.} The \emph{probability integral transform} (PIT) of the observation
$y_t$ is the predictive cumulative probability of the value actually observed:
\[
  u_t = \Pr\big(y'_t \le y_t \mid y_{1:t-1}\big)
      = \sum_{gjk}\omega_{gjk}\;
        \Phi\!\left(\frac{y_t - \mu^F_{gjk,3}}{\sqrt{\Sigma^F_{gjk,33} + R_t}}\right)
  \qquad(\code{hist\$pit}),
\]
where $y'_t$ is a draw from the predictive~\eqref{eq:pred} and $\Phi$ is the standard
normal cumulative distribution function. The plotted per-site value is the
\emph{centred} PIT of the final observation, $c_s = 2u_{T_s} - 1 \in [-1, 1]$.

\paragraph{Why this alternative.} It has the same intent and the same $[-1, 1]$ range as
$\delta_s$, but avoids its caveats:
\begin{itemize}
  \item \textbf{Out of sample:} $u_t$ uses the forecast made before $y_t$ is assimilated.
  \item \textbf{Includes observation error:} the predictive adds $R_t$.
  \item \textbf{Uses the whole distribution:} the full mixture predictive, including the
  disturbance components and the integration over process error, not a symmetric
  interval.
\end{itemize}
If observations really come from the predictive, $u_t$ is uniform on $[0, 1]$ whatever
the predictive's shape. So $c_s$ is uniform on $[-1, 1]$, and exactly 95\% of sites
should have $|c_s| \le 0.95$.

\paragraph{Panels.}
\begin{itemize}
  \item Page~2 plots $c_s$ against $y_{T_s}$, one point per site, with $[-0.95, 0.95]$
  shaded.
  \item Page~3 is the PIT histogram: all $u_t$, pooled over sites and years, as a density
  on $[0, 1]$. A calibrated filter gives a flat histogram at height 1 (dashed). The
  shapes read as:
  \begin{itemize}
    \item U-shape: the predictive is too narrow (overconfident);
    \item central hump: it is too wide;
    \item mass piled on one side: it is biased. Too many small $u_t$ means the forecast
    is too high.
  \end{itemize}
  \item Page~5 summarises each experiment: the share of sites with $|\delta_s| \le 1$,
  the median $|\delta_s|$, the share of sites with $|c_s| \le 0.95$, and the same share
  over all site-years.
\end{itemize}

%================================================================================
\section{Plot 2: probability of detection vs.\ absolute disturbance size}
\label{sec:det}

\paragraph{Detection.} Let $t_{04}$ be the timestep for 2004 ($t_{04} = 15$). The posterior
probability that site $s$ was disturbed in 2004 is
\[
  P_s = \Pr(k_{t_{04}} \ge 1 \mid y_{1:t_{04}})
      = \sum_{g}\pi_{t_{04}}(g)\sum_{j}\sum_{k\ge1} w^A_{t_{04}}(g,j,k),
\]
stored as \code{hist\$pDist} at \code{hist\$year == 2004}. It is a filtering probability:
it uses the observations up to and including 2004.

\paragraph{Size.} The x-axis is the absolute (not proportional) change in observed AGB
from 2003 to 2004:
\[
  \Delta_s = y_{s,2003} - y_{s,2004} \quad (\text{kg C m}^{-2}),
\]
which is positive for a loss. The sign is kept so that gains are visible. Points are
coloured by the product's 2004 label for the site: none, other, cut, fire or biotic.

\paragraph{Fitted curve (20 or more sites only).} When an experiment has at least 20
sites, the script fits a beta regression. This is a regression for responses in $(0,1)$
that are beta-distributed with mean $\mu_s$ and precision $\phi$:
\[
  \tilde P_s \sim \mathrm{Beta}\big(\mu_s\phi,\ (1-\mu_s)\phi\big), \qquad
  \operatorname{logit}\mu_s = \log\frac{\mu_s}{1-\mu_s} = a + b\,\Delta_s .
\]
The response is shrunk into the open interval so the beta likelihood is finite:
$\tilde P_s = \big(P_s(n-1) + 0.5\big)/n$, with $n$ the number of sites. The panel shows
the fitted mean $\mu(\Delta)$ and the band between the fitted beta distribution's 2.5\%
and 97.5\% quantiles. The summary table (page~5) lists $\Delta_{50} = -a/b$, the loss
detected with probability one half. With 5 sites no curve is fitted.

\paragraph{Reading it.} With the product, a flag raises that year's $\beta_{k,t}$ to the
user's accuracy. So a flagged site can reach a high $P_s$ even with a small $\Delta_s$.
Without the product, $P_s$ is high only when $\Delta_s$ is large compared with the spread
of the undisturbed forecast.

\paragraph{Panels.} Page~4 of \code{exps\_summary.pdf}, one panel per experiment.

%================================================================================
\section{Experiment grid and test run}
\label{sec:runs}

The \code{run-grid} chunk of \code{OR\_joint\_update2.Rmd} runs all $2\times2\times2 = 8$
experiments on the sites in \code{grid.sites}:

\begin{center}
\begin{tabular}{ll}
\toprule
factor & levels\\
\midrule
magnitude of $R_t$ & GEDI, LandTrendr\\
form of $R_t$ & heteroskedastic (\code{het}), constant (\code{con})\\
disturbance product & used (\code{TRUE}), not used (\code{FALSE})\\
\bottomrule
\end{tabular}

\smallskip
$(m_R, b_R)$: GEDI het $(0.0133, 0.0976)$, GEDI con $(0, 0.885)$,
LandTrendr het $(0.0133, 3.013)$, LandTrendr con $(0, 3.800)$.
\end{center}

All runs use the Gamma $\tau$ prior above (\code{grid.tauPrior = "gamma"}). The test run
used one randomly chosen site (seed 2004) for each of the five 2004 product labels.
Each (experiment, site) result is saved to \code{experiments/testRun} as
\begin{center}\footnotesize\path{site_<idx>_distprod_<TF>_<source>_<model>_tau<prior>.rds}\end{center}
and \path{OR_joint_update_analysis.Rmd} reads that folder (\code{indir}) and writes
\code{exps\_summary.pdf} there. For a full run, set
\code{grid.sites <- seq\_len(nrow(dextract))}.

\subsection*{Test-run results}

\begin{center}
\begin{tabular}{rlrr}
\toprule
site & 2004 product label & $\Delta_s$ (kg C m$^{-2}$) & $y_{2017}$ (kg C m$^{-2}$)\\
\midrule
11  & none   & $-0.15$ & $8.25$\\
118 & other  & $0.30$  & $2.55$\\
212 & cut    & $17.85$ & $10.80$\\
319 & fire   & $0.80$  & $2.60$\\
402 & biotic & $0.10$  & $1.90$\\
\bottomrule
\end{tabular}
\end{center}

With five sites these plots only show what each panel looks like and give a first
sanity check. They are not a calibration result.

\begin{itemize}
  \item \textbf{Filter divergence.}
  \begin{itemize}
    \item Every site is inside $[-1, 1]$ in all eight experiments. The median
    $|\delta_s|$ is $0.20$--$0.52$.
    \item It is smallest with GEDI heteroskedastic $R$ ($0.20$ and $0.23$), whose small
    $R_t$ lets the analysis follow the observation most closely.
    \item The sign follows the site: the high-biomass sites 11 and 212 have
    $\delta_s > 0$, and the three low-biomass sites have $\delta_s < 0$.
  \end{itemize}
  \item \textbf{Final-observation PIT.} Every $|c_s| \le 0.95$ as well
  ($|c_s| = 0.20$--$0.84$).
  \item \textbf{PIT histogram} (140 site-years per experiment). This is where the PIT
  shows what the divergence does not.
  \begin{itemize}
    \item The histograms are not flat. Most mass lies between $0.2$ and $0.4$.
    \item Without the product, $82$--$100\%$ of the observations at sites 11, 118, 319 and
    402 fall below the forecast median (GEDI heteroskedastic $R$). So the one-step
    forecast is biased high: VSEM grows these sites faster than observed, and the analysis
    pulls the state back each year.
    \item The share of $u_t$ in $[0.25, 0.75]$ is $0.65$--$0.78$ with LandTrendr-magnitude
    $R$, against $0.5$ expected, so that predictive is also too wide. With GEDI
    heteroskedastic $R$ it is $0.37$--$0.49$.
    \item With the product, a second peak appears near $u = 0.8$. It comes from flagged
    years with little observed loss: the product's flag puts predictive mass on low
    biomass, so an unchanged observation lands in the upper tail (e.g.\ site 118 in 2004:
    $u = 0.82$ with the product, $0.32$ without).
  \end{itemize}
  \item \textbf{Detection.}
  \begin{itemize}
    \item Site 212 has $P_s = 1.000$ in all eight experiments: its loss of $17.85$ is
    detected from the biomass data alone.
    \item The other sites lose at most $0.8$. Without the product they have
    $P_s \le 0.061$.
    \item With the product, the flagged sites reach $0.36$--$0.83$, and the unflagged
    site 11 stays at $0.015$--$0.035$.
  \end{itemize}
\end{itemize}

\begin{figure}[p]
  \centering
  \includegraphics[width=\textwidth,page=1]{exps_summary.pdf}
  \caption{Filter divergence $\delta_s$ at the final observation, one point per site and
           one panel per experiment (test run).}
  \includegraphics[width=\textwidth,page=2]{exps_summary.pdf}
  \caption{Centred PIT $c_s$ of the final observation, one point per site (test run).}
\end{figure}
\begin{figure}[p]
  \centering
  \includegraphics[width=\textwidth,page=3]{exps_summary.pdf}
  \caption{PIT histograms over all sites and years (test run).}
  \includegraphics[width=\textwidth,page=4]{exps_summary.pdf}
  \caption{Posterior probability of disturbance in 2004 against the observed AGB loss
           from 2003 to 2004 (test run).}
\end{figure}

\end{document}
